% BHU PYQ -- Professional Exam Template
\documentclass[11pt,a4paper]{article}

\usepackage[a4paper,top=2.5cm,bottom=2.5cm,left=2.8cm,right=2.8cm,headheight=14pt]{geometry}
\usepackage{amsmath,amssymb}
\usepackage[T1]{fontenc}
\usepackage[utf8]{inputenc}
\usepackage{lmodern}
\usepackage{microtype}
\usepackage{fancyhdr}
\usepackage{enumitem}
\usepackage{listings}
\usepackage{hyperref}
\usepackage{lastpage}
\usepackage{parskip}

\hypersetup{colorlinks=true,urlcolor=black,linkcolor=black,
  pdftitle={CSB-503: Discrete Mathematical Structure -- BHU PYQ}}

\pagestyle{fancy}
\fancyhf{}
\renewcommand{\headrulewidth}{0.4pt}
\renewcommand{\footrulewidth}{0.4pt}
\fancyhead[L]{\small\textit{Banaras Hindu University}}
\fancyhead[R]{\small\textit{CSB-503: Discrete Mathematical Structur}}
\fancyfoot[C]{\small Page \thepage\ of \pageref{LastPage}}

\lstset{
  basicstyle=\small\ttfamily,
  frame=single,
  breaklines=true,
  columns=flexible,
  keepspaces=true
}

\newlist{parts}{enumerate}{2}
\setlist[parts,1]{label=(\alph*),leftmargin=*,itemsep=4pt,topsep=3pt}
\setlist[parts,2]{label=(\roman*),leftmargin=*,itemsep=2pt,topsep=2pt}
\newcommand{\pts}[1]{\hfill\textit{[#1]}}

\begin{document}
\begin{center}
  {\large\bfseries BANARAS HINDU UNIVERSITY}\\[2pt]
  {\normalsize B.Sc. (Hons.) Semester V Examination, 2013-14}\\[6pt]
  \rule{0.8\linewidth}{0.5pt}\\[6pt]
  {\large\bfseries Paper: CSB-503 --- Discrete Mathematical Structure}\\[4pt]
  {\normalsize Time: Three hours \hfill Maximum Marks: 70}
\end{center}

\rule{\linewidth}{0.4pt}

\smallskip
\noindent\textit{Instructions: Attempt any FIVE questions, The figures in right margin indicate the marks.}

\bigskip

\medskip\noindent\textbf{Question 1.}
\begin{parts}
  \item Define successor of a set. Give two sets $A$ and $B$ such that $A \in B$ and $A \subset B$. \pts{4}
  \item Define difference and symmetric difference of sets. Describe in general how these two set operations are achieved when sets are represented using bit strings. \pts{5}
  \item Define partition and covering of a set. Generate universal set $E$ using $A_1 = \{3, 4, 5\}$, $A_2 = \{1, 2\}$, and then give the partition of generated universal set using sets $A_1, A_2$. \pts{5}
\end{parts}

\medskip\noindent\textbf{Question 2.}
\begin{parts}
  \item Consider the following relation $R$ on set $A = \{1, 2, 3, 4\}$:
  \[ R = \{(a, b) \mid a \text{ divides } b\} \]
  \begin{parts}
    \item List all ordered pairs of relation $R$. \pts{1}
    \item Represent relation $R$ using digraph. \pts{1}
    \item Give relation matrix of the relation $R$. \pts{1}
    \item Find the matrix of $R^{-1}$ from matrix of relation $R$. \pts{1}
  \end{parts}
  \item 
  \begin{parts}
    \item Define partial order relation and poset. \pts{2}
    \item Using only matrix operations show that the relation $R$ defined in Question 2(a) is a partial order relation and hence $(A, R)$ is a poset. \pts{3}
    \item Draw the Hasse diagram of the above partial order relation. \pts{1}
    \item Find the maximal and the minimal elements, if they exist, of the above poset $(A, R)$. \pts{2}
    \item Find the greatest and the least elements, if they exist, of the above poset $(A, R)$. \pts{2}
  \end{parts}
\end{parts}

\medskip\noindent\textbf{Question 3.}
\begin{parts}
  \item Let $A$ be a set having $n$ elements. Show the following:
  \begin{parts}
    \item Total possible reflexive relations on $A$ is $2^{n(n-1)}$. \pts{2}
    \item Total possible symmetric relations on $A$ is $2^{n(n+1)/2}$. \pts{2}
    \item Total possible symmetric and reflexive relations on $A$ is $2^{n(n-1)/2}$. \pts{2}
    \item Total possible anti-symmetric relations on $A$ is $2^n \cdot 3^{n(n-1)/2}$. \pts{2}
  \end{parts}
  \item Define number theoretic function, partial function, and primitive recursive function and give at least one example of each. \pts{4}
  \item Show that the total number of one-to-one functions from $A$ to $B$ is $n(n-1)(n-2)\cdots(n-m+1)$, where $|A| = m$ and $|B| = n$. \pts{2}
\end{parts}

\medskip\noindent\textbf{Question 4.}
\begin{parts}
  \item Prove the following without constructing truth table:
  \begin{parts}
    \item $\neg p \land (p \lor q) \to q$ \pts{3}
    \item $[p \land (p \to q)] \to q$ \pts{3}
  \end{parts}
  \item Represent the following facts using propositions. Also find what rules of inference are used in each of these arguments: \pts{6}
  \begin{parts}
    \item Alice is a mathematics major. Therefore, Alice is either a mathematics major or a computer science major.
    \item Jerry is a mathematics major and a computer science major. Therefore, Jerry is a mathematics major.
    \item If I go swimming, then I will stay in the sun too long. If I stay in the sun too long, then I will get sunburned. Therefore, if I go swimming, then I will get sunburned.
  \end{parts}
  \item What is the limitation of statement calculus and how is it removed in predicate calculus? \pts{2}
\end{parts}

\medskip\noindent\textbf{Question 5.}
\begin{parts}
  \item 
  \begin{parts}
    \item Prove that in a group $(G, \circ)$, if every element is its own inverse, i.e., $a^2 = e$ for all $a \in G$, then $G$ is an abelian group. \pts{5}
    \item Let $(\mathbb{R}, +)$ denote the additive group of real numbers and $(\mathbb{R}^*, \cdot)$ denote the multiplicative group of positive reals. Define a function $f: \mathbb{R} \to \mathbb{R}^*$ as $f(x) = a^x$, $a > 0$. Show that $f$ is a homomorphism. \pts{5}
  \end{parts}
  \item Let $f$ be a homomorphism from $(\mathbb{Z}, +)$ to $(\mathbb{Q}^* , \cdot)$, $\mathbb{Q}^* = \mathbb{Q} - \{0\}$. Suppose $f(2) = 4$. Find out $f(6)$ and $f(-8)$. \pts{4}
\end{parts}

\medskip\noindent\textbf{Question 6.}
\begin{parts}
  \item Model the round-robin tournament using a graph. \pts{3}
  \item What do the in-degree and out-degree of a vertex in the graph represent? \pts{3}
  \item Draw the graph of the following: \pts{4}
  \begin{parts}
    \item $K_5$
    \item $C_5$
    \item $W_5$
    \item $Q_3$
  \end{parts}
  \item For which values of $n$ are these graphs bipartite? Also find their chromatic number: \pts{4}
  \begin{parts}
    \item $K_n$
    \item $C_n$
    \item $W_n$
    \item $Q_n$
  \end{parts}
\end{parts}

\medskip\noindent\textbf{Question 7.}
\begin{parts}
  \item Can five houses be connected to two utilities without connections crossing? \pts{2}
  \item A connected planar graph has 30 edges. If a planar representation of this graph divides the plane into 20 regions, how many vertices does this graph have? \pts{4}
  \item Find the complements of every element, if they exist, of $(S_n, D)$ for $n = 75$, where $S_n$ is the set of divisors of $n$, and $D$ is the division relation. \pts{5}
  \item Determine whether the following statements are true, false, or not valid: \pts{3}
  \begin{parts}
    \item $2+2=4$ iff $1+1=2$.
    \item If $1+1=3$, then unicorns exist.
    \item $x + y = 2$.
  \end{parts}
\end{parts}

\vfill
\rule{\linewidth}{0.4pt}
\begin{center}\small
\href{https://bhu-exam.vercel.app/}{Click here to visit our website}\\[4pt]\href{https://me-aryan.vercel.app/}{Click here to visit our contributor}\\[4pt]\href{https://quashan-me.vercel.app/}{Click here to visit our contributor}
\end{center}
\end{document}